\documentclass{beamer}
\usetheme{Madrid}
\usecolortheme{default}

\usepackage{amsmath}
\usepackage{graphicx}

\title[Smart Parking System]{An Enhanced MARL Framework for Smart Parking \\ with Route Optimization and Environmental Impact Analysis}
\author[P. Simon Peter, M. Dinesh Naik]{1\textsuperscript{st} P. Simon Peter \and 2\textsuperscript{nd} M. Dinesh Naik}
\institute[IIIT Sri City]{
  Computer Science and Engineering \\
  IIIT Sri City
}
\date{2026}

\begin{document}

% Page 1: Title Slide
\begin{frame}
  \titlepage
\end{frame}

% Page 2: Introduction & Problem Statement
\begin{frame}{Introduction}
  \begin{itemize}
      \item \textbf{The Problem:} Finding available parking in urban areas contributes to traffic congestion, driver frustration, and environmental degradation.
      \item \textbf{Traditional Methods:} Methods like First-In-First-Out (FIFO) or Service-In-Random-Order (SIRO) fail to adapt to dynamic urban environments.
      \item \textbf{The Need:} A smart, adaptive system that can intelligently allocate parking spots and guide users effectively.
  \end{itemize}
\end{frame}

% Page 3: Baseline Methodology (AI Models)
\begin{frame}{Baseline Methodology}
  The foundation relies on an AI-driven predictive and allocative framework:
  \vspace{0.5cm}
  \begin{itemize}
      \item \textbf{Multi-Agent Reinforcement Learning (MARL):} Models competitive interactions among vehicles seeking parking, maximizing rewards based on parking success.
      \item \textbf{Graph Neural Networks (GNN):} Captures complex spatial dependencies and the topological structure of the parking network.
      \item \textbf{Gated Recurrent Units (GRU):} Processes historical data to capture temporal dynamics and predict future demand.
  \end{itemize}
\end{frame}

% Page 4: Limitations of Existing Work
\begin{frame}{Limitations of Existing Work}
  Despite high accuracy in predicting and allocating spots, state-of-the-art MARL-GNN-GRU frameworks have critical gaps:
  \vspace{0.5cm}
  \begin{itemize}
      \item \textbf{Lack of Practical Guidance:} They allocate a spot theoretically but don't optimize the physical journey for the user.
      \item \textbf{No User Interface:} They lack user-facing applications (mobile apps) for real-world deployment.
      \item \textbf{Unquantified Sustainability:} The intuitive environmental benefits are widely acknowledged but rarely quantified in a concrete manner.
  \end{itemize}
\end{frame}

% Page 5: Proposed Enhancements Overview
\begin{frame}{Proposed Enhancements Overview}
  To bridge the gap between theory and practice, we introduce three novel contributions:
  \vspace{0.5cm}
  \begin{enumerate}
      \item \textbf{Route Optimization Module:} A multi-criteria routing system minimizing total commute time.
      \item \textbf{Mobile Application Integration:} A robust, user-centric mobile app built with React Native.
      \item \textbf{Environmental Impact Calculator:} A quantifiable metric system to track CO$_2$ emissions saved per search.
  \end{enumerate}
\end{frame}

% Page 6: Route Optimization Module
\begin{frame}{Route Optimization Module}
  Minimizes commute by calculating the distance between the user, parking spot, and destination using the Haversine formula:
  \begin{equation*}
      d = 2r \arcsin\left(\sqrt{\sin^2\left(\frac{\Delta \phi}{2}\right) + \cos\phi_1\cos\phi_2\sin^2\left(\frac{\Delta \lambda}{2}\right)}\right)
  \end{equation*}
  
  \textbf{Where:}
  \begin{itemize}
      \item $d$: Distance \quad $r$: Earth's radius
      \item $\phi_1, \phi_2$: Latitudes \quad $\Delta \phi, \Delta \lambda$: Differences in latitude and longitude
  \end{itemize}
  
  \textbf{Multi-Criteria Scoring} evaluates spots based on:
  \begin{itemize}
      \item \textbf{Drive Time:} Time to reach the parking spot.
      \item \textbf{Walk Time:} Time to walk to the final destination.
      \item \textbf{Availability Score:} MARL-predicted likelihood of the spot remaining free.
  \end{itemize}
\end{frame}

% Page 7: Mobile Application (React Native)
\begin{frame}{Mobile Application (React Native)}
  Transitioning the theoretical AI model into a usable product:
  \vspace{0.5cm}
  \begin{itemize}
      \item \textbf{Cross-Platform Framework:} Developed using React Native (with Expo) for seamless iOS and Android support.
      \item \textbf{Real-Time Map Interface:} Integrates MapLibre GL JS via WebView to visualize the parking landscape.
      \item \textbf{Live Availability Indicators:} Spots are color-coded (Green = Available, Red = Occupied) using real-time GRU predictions.
      \item \textbf{One-Touch Navigation:} A ``Find Parking'' feature instantly triggers route optimization and turn-by-turn navigation.
  \end{itemize}
\end{frame}

% Page 8: Environmental Impact Calculator
\begin{frame}{Environmental Impact Calculator}
  Quantifies sustainability benefits by calculating the reduction in CO$_2$ emissions ($\Delta \text{CO}_2$):
  \begin{equation*}
      \Delta \text{CO}_2 = (v \times t_{\text{cruise}} - D_{\text{smart}}) \times E_f
  \end{equation*}
  
  \textbf{Where:}
  \begin{itemize}
      \item $v$: Average cruising speed (e.g., 20 km/h)
      \item $t_{\text{cruise}}$: Traditional cruising time (e.g., 8 mins)
      \item $D_{\text{smart}}$: Optimized route distance (e.g., 0.5 km)
      \item $E_f$: Vehicle emission factor (e.g., 120g CO$_2$/km)
  \end{itemize}
  
  \vspace{0.2cm}
  \textit{Result:} A net saving of \textbf{260g of CO$_2$} per search, tracked cumulatively in the app to promote eco-friendly behavior.
\end{frame}

% Page 9: Impact and Discussion
\begin{frame}{Impact and Discussion}
  \begin{itemize}
      \item \textbf{Reduction in Search Time:} Direct navigation reduces actual search time by 30-40\%, alleviating localized congestion.
      \item \textbf{Practical Usability:} The intuitive Flutter UI ensures sophisticated AI predictions are easily accessible.
      \item \textbf{Sustainability Tracking:} Gamifies eco-conscious behavior for users and provides city planners with tangible environmental efficacy data.
  \end{itemize}
\end{frame}

% Page 10: Conclusion
\begin{frame}{Conclusion \& Future Work}
  \begin{itemize}
      \item \textbf{Summary:} We successfully transformed a theoretical MARL-GNN-GRU resource allocation model into a practical, sustainable urban mobility solution.
      \item \textbf{Significance:} The integration of Route Optimization, Mobile App, and Environmental tracking bridges the gap between backend AI and end-user needs.
      \item \textbf{Future Work:} Deploying the system in a live urban environment and integrating real-time traffic data to further refine routing algorithms.
  \end{itemize}
  \vspace{1cm}
  \begin{center}
      \textbf{Thank You!}
  \end{center}
\end{frame}

\end{document}
